\documentclass[a4paper,11pt]{article}

\usepackage[spanish]{babel}
\usepackage[T1]{fontenc}
\usepackage[hmargin=0.65in, vmargin=0.5in]{geometry}
\usepackage{tabularx}
\usepackage{xcolor}
\usepackage{titlesec}
\usepackage{enumitem}
\usepackage{hyperref}

\definecolor{accent}{HTML}{0F4C5C}
\definecolor{muted}{HTML}{5F6B73}
\hypersetup{colorlinks=true, urlcolor=accent, linkcolor=accent}

\renewcommand{\familydefault}{\sfdefault}
\linespread{1.05}
\pagestyle{empty}
\setlength{\parindent}{0pt}
\raggedright

\titleformat{\section}{\Large\bfseries\color{accent}}{}{0em}{}[{\color{accent}\titlerule[0.8pt]}]
\titlespacing*{\section}{0pt}{11pt}{5pt}

\setlist[itemize]{label={\color{accent}\textbullet}, leftmargin=1.4em, topsep=2pt, itemsep=2pt, parsep=0pt}

% \entry{Organización}{Fecha}{Cargo}
\newcommand{\entry}[3]{%
  {\large\textbf{#1}} \hfill {\color{muted}#2} \\
  {\color{accent}\textit{#3}}\par
}
\newcommand{\entrygap}{\vspace{7pt}}

\begin{document}

% Encabezado
\begin{center}
  {\huge\bfseries Alberto Zúñiga Marinovic} \\[5pt]
  {\large\color{accent} Ingeniero Civil en Ciencias de la Computación} \\[2pt]
  {\color{muted} Universidad de los Andes} \\[5pt]
  +56 9 6496 2736 \quad$\cdot$\quad
  \href{mailto:a.zuniga.marinovic@gmail.com}{a.zuniga.marinovic@gmail.com} \quad$\cdot$\quad
  \href{https://www.linkedin.com/in/alberto-zuniga-marinovic}{LinkedIn} \quad$\cdot$\quad
  \href{https://github.com/AlbertoZuiga}{GitHub} \quad$\cdot$\quad
  \href{https://albertozuiga.github.io}{Portafolio}
\end{center}

\section*{Experiencia}

\entry{Buk}{Mar 2026 -- Presente}{Software Engineer Level 1, equipo de Talento - Evaluaciones}
\begin{itemize}
  \item Desarrollo y mantenimiento del módulo de \textbf{Evaluaciones de Desempeño} de la plataforma.
\end{itemize}
\entrygap

\entry{Buk}{Ene 2026 -- Feb 2026}{Práctica profesional como Software Engineer, equipo de Talento - Evaluaciones}
\begin{itemize}
  \item Migración de la sección \textbf{``Mis evaluaciones''} a \textbf{Fiji}, el framework de desarrollo interno de Buk.
\end{itemize}
\entrygap

\entry{Universidad de los Andes}{Ago 2022 -- Jun 2026}{Ayudante}
\begin{itemize}
  \item \textbf{14 ayudantías en 9 cursos:} Paradigmas de Programación (×4), Web Technologies (×3), Programación (×2),
        Bases de Datos, Aplicaciones Móviles, Sistemas Electrónicos, Taller de Computación, Taller de Proyectos de Ingeniería.
\end{itemize}

\section*{Proyectos}


{\large\textbf{Scheduler App}} \quad {\color{accent}\textit{Python, Flask}} \hfill \href{http://scheduler-app-iu34.onrender.com/}{Ver demo}
\begin{itemize}
  \item Gestión de horarios y división automática de grupos.
\end{itemize}
\entrygap

{\large\textbf{Healthy}} \quad {\color{accent}\textit{Ruby on Rails}} \hfill \href{https://healthy-k6hn.onrender.com}{Ver demo}
\begin{itemize}
  \item Sistema de recomendación de planes de comida saludables integrado con compras.
\end{itemize}

\section*{Educación}

\entry{Universidad de los Andes}{2020 -- 2026}{Ingeniero Civil en Ciencias de la Computación}
\begin{itemize}
  \item Concentración tecnológica en Ingeniería Civil Eléctrica. Minor en Psicología.
\end{itemize}

\section*{Liderazgo y voluntariado}

\entry{Pastoral Universidad de los Andes}{2025 -- 2026}{Catequista (2025 -- 2026) y encargado de logística (2026)}
\begin{itemize}
  \item Gestión logística y organización de las actividades del grupo.
  \item Planificación y dictado de clases semanales de formación.
\end{itemize}
\entrygap

\textbf{Consejero político de Ingeniería Civil:} representante ante la Federación de Estudiantes. \hfill {\color{muted}2023} \\[4pt]
\textbf{Fundación Nueva Mente:} voluntario pro bono; sitio web (Wix) y registro de gastos. \hfill {\color{muted}2022 -- 2024}

\section*{Habilidades}

{\renewcommand{\arraystretch}{1.15}
\begin{tabularx}{\textwidth}{@{}>{\bfseries\color{accent}}l X@{}}
  Lenguajes    & Python, JavaScript, Ruby, SQL, C++, HTML/CSS \\
  Frameworks   & React, React Native, Flask, Django, FastAPI \\
  Herramientas & Git/GitHub, Docker, LaTeX, Excel \\
  Idiomas      & Español (nativo), Inglés (intermedio) \\
\end{tabularx}}

\end{document}
